% ---------------------------------------------------------------------------
% Front matter: title page, about this compendium, legend, provenance index
% ---------------------------------------------------------------------------
\begin{titlepage}
\centering
\vspace*{1.6in}
{\Huge\bfseries New Mathematics for\\[4pt] Marketing Measurement\par}
\vspace{14pt}
{\LARGE A Compendium of the Deep-Dive Program\par}
\vspace{10pt}
{\large Runs 01--22, 31 August -- 7 September 2026\par}
\vspace{1.4in}
{\large Compiled and reviewed 7 September 2026\par}
\vspace{6pt}
{\normalsize from \texttt{marketing\_measurement/frontier\_study/deep\_dives/}\par}
\vfill
\begin{minipage}{0.86\textwidth}\small
This document stitches together twenty-two automated research runs into fifteen self-contained chapters.
Each chapter reproduces the background ideas its dive builds on, states the dive's constructions and results
with the dive's own epistemic labels, tabulates the headline numerical evidence, and closes with a
scientific-paper style discussion of impacts and further exploration. Compiler's reviewer notes (orange boxes)
flag what a referee would want checked before believing a claim. Appendix~\ref{app:ledger} is the ledger of all
88 open problems the program spawned and their status at compilation.
\end{minipage}
\vspace{0.5in}
\end{titlepage}

\chapter*{About this compendium}
\addcontentsline{toc}{chapter}{About this compendium}
\markboth{About this compendium}{}

\section*{What was compiled}

Between 31 August and 7 September 2026 a scheduled task executed twenty-two ``deep dives'' into the open
problems catalogued by the frontier study of marketing measurement (\texttt{01\_state\_of\_the\_field.md},
\texttt{02\_open\_questions.md}, \texttt{03\_mmm\_adoption\_barriers.md}). Each run took one item from a queue,
worked it under a fixed depth protocol (formalisation with explicit assumptions, a literature search, at least
three construct--attack--refine rounds, at least eight numerical experiments including robustness sweeps and a
negative or power control, a clean-room re-derivation of at least one core result by an independent subagent,
uncertainty on every headline number, and a devil's-advocate paragraph against every headline claim), and wrote
a markdown report with accompanying code. Every fourth run was an \emph{extension} that pushed the most
promising prior dive further. The program's index file (\texttt{deep\_dives/00\_INDEX.md}) holds the queue,
the backlog of open problems each dive spawned (``BL'' items), and a one-line log of each run's finding.

This compendium was assembled on 7 September 2026 from the twenty-two reports and their code as they stood on
disk at that date (the provenance index below records each file's last-written timestamp). Its purpose is
to let a reader review, on paper, what ideas were actually developed---without having to open the dives, the
papers they cite, or the index.

\section*{How the chapters are organised}

There is one chapter per dive, except that a dive and its extension are merged into a single chapter
(dives 03+04, 08+09, 13+14, 16+19), and the four dives that form one continuous thread on identification
(01, then the Phase D chain 20, 21, 22) are merged into \cref{ch:identification}. Chapters otherwise follow
the program's run order. Every chapter has the same skeleton:

\begin{description}[leftmargin=1.6em,itemsep=1pt]
\item[Provenance] which dives, files, code and frontier-study problem the chapter covers, and which backlog items it opened or closed.
\item[Key ideas] the chapter's central results in a few sentences, with their headline numbers.
\item[Problem and motivation] what the dive set out to answer and what the field believed before.
\item[Background] the ideas the dive builds on, reproduced in summary so that the chapter is self-contained: for every paper or established result the dive synthesises, its central construction or equation, what it assumes, and what it delivers.
\item[Setup and assumptions] the formal model and notation.
\item[Main results] the dive's constructions, theorems and laws, with key derivation steps and the dive's own epistemic labels (see the legend).
\item[Numerical evidence] the headline experiments as tables, with the uncertainty the dive reports, including negative and power controls.
\item[Limitations and the devil's advocate] the dive's own strongest counterarguments and what the attack-and-refine rounds changed.
\item[Discussion: impacts] what changes for practice and for theory.
\item[Further exploration] the dive's next steps, the backlog items it opened, and the compiler's view of the most promising extension.
\item[Reviewer's notes] the compiler's own flags, as a referee would write them.
\item[Sources] the papers summarised or relied on, with links where the program's bibliographies had them.
\end{description}

\section*{Legend}

Results carry the label the dive itself gave them. \tagEstablished\ marks a result derived in the dive and, where
the dive says so, reproduced by its clean-room verifier. \tagTransported\ marks a known result carried into the
MMM setting from another literature. \tagNumerical\ marks a result established by simulation rather than proof.
\tagConjecture\ marks speculation the dive labelled as such. \tagRefuted\ marks a claim the dive itself made and
then overturned in a later round---these are kept because the refutations are among the more instructive
findings. \dive{NN} refers to run NN; \BL{NN} to backlog item NN in \cref{app:ledger}. Orange boxes are the
compiler's reviewer notes and are the only content in the chapters that does not come from the dives.

Two program-wide caveats apply to every number in this book. First, all results are from simulation on
synthetic data under the model each dive states; none has been run on an advertiser's data. Second,
\dive{21}'s verifier found that pseudo-inverses at default tolerance had silently manufactured finite standard
errors from singular Fisher matrices; \dive{22} audited dives 20--22 (\BL{82}) but the FIM-based standard
errors in earlier dives (notably 02, 09, 16, 19) have not had that audit. The reviewer notes say where this
matters.

\section*{Provenance index}

\Cref{tab:provenance} lists every source dive, its file, the run date recorded in the program log, the
file's last-written timestamp on disk at compilation, its length, and the chapter that covers it. Each dive's
code lives beside it in \texttt{deep\_dives/code/} under the same stem.

\begin{table}[htbp]
\centering\small
\caption{Source dives and their dates (UTC). ``Run'' is the start time recorded in the program log; ``Written'' is the report file's last modification time on disk on 2026-09-07. Word counts are of the markdown reports.}
\label{tab:provenance}
\begin{tabular}{@{}l l l l r c@{}}
\toprule
Dive & File (\texttt{deep\_dives/}) & Run (log) & Written & Words & Ch. \\
\midrule
01 & \texttt{01\_mmm\_identified\_set.md} & 2026-08-31 19:00 & 2026-08-31 22:08 & 2{,}003 & 1 \\
02 & \texttt{02\_experiment\_transport\_map.md} & 2026-08-31 21:15 & 2026-08-31 23:14 & 2{,}225 & 2 \\
03 & \texttt{03\_decision\_voi.md} & 2026-09-01 00:21 & 2026-09-01 00:21 & 1{,}582 & 3 \\
04 & \texttt{04\_decision\_voi\_ext.md} (ext.\ of 03) & 2026-09-01 02:30 & 2026-09-01 01:25 & 1{,}982 & 3 \\
05 & \texttt{05\_fusion\_loop.md} & 2026-09-01 05:30 & 2026-09-01 03:39 & 2{,}184 & 4 \\
06 & \texttt{06\_aggregation\_privacy.md} & 2026-09-01 07:45 & 2026-09-01 07:36 & 2{,}582 & 5 \\
07 & \texttt{07\_refresh\_stability.md} & 2026-09-01 12:30 & 2026-09-01 11:56 & 2{,}507 & 6 \\
08 & \texttt{08\_decision\_deadband.md} & 2026-09-01 15:30 & 2026-09-01 16:55 & 2{,}936 & 7 \\
09 & \texttt{09\_decision\_deadband\_ext.md} (ext.\ of 08) & 2026-09-01 20:05 & 2026-09-01 20:01 & 3{,}933 & 7 \\
10 & \texttt{10\_algorithmic\_delivery\_estimands.md} & 2026-09-01 23:30 & 2026-09-01 23:43 & 3{,}962 & 8 \\
11 & \texttt{11\_data\_archaeology\_compiler.md} & 2026-09-01 21:20\,$^{a}$ & 2026-09-02 03:33 & 4{,}024 & 9 \\
12 & \texttt{12\_surrogacy\_horizon\_bounds.md} & 2026-09-02 04:50 & 2026-09-02 11:34 & 3{,}640 & 10 \\
13 & \texttt{13\_attribution\_impossibility.md} & 2026-09-02 15:40 & 2026-09-02 15:40 & 4{,}535 & 11 \\
14 & \texttt{14\_attribution\_impossibility\_ext.md} (ext.\ of 13) & 2026-09-02 19:45 & 2026-09-02 19:44 & 4{,}536 & 11 \\
15 & \texttt{15\_calibrated\_simulator.md} & 2026-09-02 23:30 & 2026-09-02 23:32 & 4{,}258 & 12 \\
16 & \texttt{16\_feedback\_loop\_inference.md} & 2026-09-03 03:30 & 2026-09-03 03:44 & 5{,}149 & 13 \\
17 & \texttt{17\_human\_mechanism\_panel.md} & 2026-09-03 10:30 & 2026-09-03 14:24 & 5{,}021 & 14 \\
18 & \texttt{18\_heavy\_tail\_efficiency.md} & 2026-09-03 16:10 & 2026-09-03 16:01 & 5{,}642 & 15 \\
19 & \texttt{19\_feedback\_loop\_inference\_ext.md} (ext.\ of 16) & 2026-09-03 14:45\,$^{a}$ & 2026-09-03 21:41 & 5{,}489 & 13 \\
20 & \texttt{20\_price\_of\_identification.md} & 2026-09-04 09:55 & 2026-09-04 17:20 & 5{,}236 & 1 \\
21 & \texttt{21\_spectral\_comb.md} & 2026-09-06 11:20 & 2026-09-06 19:17 & 6{,}293 & 1 \\
22 & \texttt{22\_price\_of\_the\_prior.md} & 2026-09-07 16:45 & 2026-09-07 16:46 & 6{,}484 & 1 \\
\midrule
\multicolumn{6}{@{}l@{}}{\footnotesize $^{a}$ The log time is earlier than the file's own timestamps and than the preceding run's; it is taken as a logging error.} \\
\bottomrule
\end{tabular}
\end{table}

The program's index file was last written at 2026-09-07 16:47 UTC and lists run 23 (an extension) as the
next item in the queue; nothing from run 23 or later is in this compendium.
